\documentclass[11pt]{beamer}

% ============================================================
% Theme & Appearance
% ============================================================
\usetheme{Madrid}
\usecolortheme{seahorse}
\setbeamertemplate{navigation symbols}{}

% ============================================================
% Encoding & Fonts (pdfLaTeX-safe)
% ============================================================
\usepackage[utf8]{inputenc}
\usepackage[T1]{fontenc}
\usepackage{lmodern}

% ============================================================
% Math & Graphics
% ============================================================
\usepackage{amsmath, amssymb}
\usepackage{graphicx}
\usepackage{tikz}
\usetikzlibrary{arrows.meta}

% ============================================================
% Code Listings
% ============================================================
\usepackage{listings}
\usepackage{xcolor}

\lstset{
	basicstyle=\ttfamily\small,
	frame=single,
	breaklines=true,
	columns=fullflexible,
	keywordstyle=\color{blue},
	commentstyle=\color{gray},
	stringstyle=\color{red},
	showstringspaces=false
}

% ============================================================
% Title Information
% ============================================================
\title{Concurrency \& Distributed Systems}
\subtitle{Week 1 --- Part 3: Theoretical Framing (Why Concurrency Works and Fails)}
\author{Dr. Mohamad Aoude}
\institute{Lebanese University \\ Faculty of Engineering}
\date{\today}

% ============================================================
\begin{document}
	
	% ============================================================
	\begin{frame}
		\titlepage
	\end{frame}
	
	% ============================================================
	\section{Theoretical Framing}
	% ============================================================
	
	% ------------------------------------------------------------
	\begin{frame}{Why Concurrency Works (and Fails)}
		Practical systems fail for mathematical reasons.
		

		Before adding threads blindly, we must understand three principles:
		
		\begin{itemize}
			\item \textbf{Amdahl's Law} (limits of speedup)
			\item \textbf{Little's Law} (queue behavior)
			\item \textbf{Context switching cost} (hidden overhead)
		\end{itemize}
		

		\begin{alertblock}{Senior-level mindset}
			You must justify performance claims with math \textbf{and} measurements.
		\end{alertblock}
	\end{frame}
	
	% ============================================================
	\section{Amdahl's Law}
	% ============================================================
	
	% ------------------------------------------------------------
	\begin{frame}{Amdahl's Law --- The Ceiling}
		\[
		S(N) = \frac{1}{(1 - P) + \frac{P}{N}}
		\]
		

		Where:
		\begin{itemize}
			\item $P$ = parallelizable fraction
			\item $N$ = number of cores/threads
			\item $S(N)$ = maximum theoretical speedup
		\end{itemize}
		

		\begin{block}{Interpretation}
			The sequential part $(1-P)$ dominates as $N$ grows.
		\end{block}
	\end{frame}
	
	% ------------------------------------------------------------
	\begin{frame}{Worked Example}
		Assume:
		\begin{itemize}
			\item $P = 0.8$ (80\% parallelizable)
			\item $N = 10$
		\end{itemize}
		

		\[
		S(10) = \frac{1}{0.2 + \frac{0.8}{10}}
		= \frac{1}{0.2 + 0.08}
		= \frac{1}{0.28}
		\approx 3.57
		\]

		\begin{alertblock}{Key point}
			Even 10 threads cannot give 10$\times$ speedup.
		\end{alertblock}
	\end{frame}
	
	% ------------------------------------------------------------
	\begin{frame}{Infinite Cores Limit}
		As $N \to \infty$:
		\[
		S_{\max} = \frac{1}{1-P}
		\]
		

		If $P=0.8$:
		\[
		S_{\max} = \frac{1}{0.2} = 5
		\]
		

		\begin{block}{Meaning}
			Even infinite cores cannot exceed 5$\times$ speedup.
		\end{block}
	\end{frame}
	
	% ------------------------------------------------------------
	\begin{frame}{Student Calculation (1 minute)}
		\textbf{Question:} If $P=0.5$, what is the maximum speedup with infinite cores?
		

		\[
		S_{\max} = \frac{1}{1-0.5} = 2
		\]
		

		\begin{alertblock}{Lesson}
			If half the program is sequential, scaling is fundamentally limited.
		\end{alertblock}
	\end{frame}
	
	% ============================================================
	\section{Live Demo: AmdahlCalculator}
	% ============================================================
	
	% ------------------------------------------------------------
	\begin{frame}[fragile]{Student Instructions: Run AmdahlCalculator (Mandatory)}
		\textbf{Step 1 --- Compile}
		\begin{lstlisting}
			javac AmdahlCalculator.java
		\end{lstlisting}
		

		\textbf{Step 2 --- Run}
		\begin{lstlisting}
			java AmdahlCalculator
		\end{lstlisting}
		

		\textbf{Step 3 --- Record}
		\begin{itemize}
			\item Copy the printed table into your report (or take a screenshot)
			\item Highlight the ``inf cores'' column
		\end{itemize}
	\end{frame}
	
	% ------------------------------------------------------------
	\begin{frame}[fragile]{AmdahlCalculator.java (Core Function)}
		Focus on the formula:
		
		\begin{lstlisting}[language=Java]
			public static double calculateSpeedup(double P, int N) {
				return 1.0 / ((1.0 - P) + (P / N));
			}
		\end{lstlisting}
		

		\begin{block}{What to watch}
			As $N$ grows, the speedup increases then \textbf{plateaus}.
		\end{block}
	\end{frame}
	
	% ------------------------------------------------------------
	\begin{frame}{Expected Pattern (Sanity Check)}
		Your exact rounding may differ, but the \textbf{limits} must match:
		
		\begin{itemize}
			\item $P=0.50 \Rightarrow S_{\max} \approx 2.00$
			\item $P=0.75 \Rightarrow S_{\max} \approx 4.00$
			\item $P=0.90 \Rightarrow S_{\max} \approx 10.00$
			\item $P=0.95 \Rightarrow S_{\max} \approx 20.00$
		\end{itemize}
		

		\begin{alertblock}{If your output does not match}
			Your formula implementation is wrong.
		\end{alertblock}
	\end{frame}
	
	% ============================================================
	\section{Little's Law}
	% ============================================================
	
	% ------------------------------------------------------------
	\begin{frame}{Little's Law --- Why Queues Form}
		\[
		L = \lambda W
		\]
		

		Where:
		\begin{itemize}
			\item $L$ = average number of requests in the system
			\item $\lambda$ = arrival rate (req/s)
			\item $W$ = average time in system (seconds)
		\end{itemize}
		

		\begin{block}{Key idea}
			If latency $W$ increases, queue size $L$ increases automatically.
		\end{block}
	\end{frame}
	
	% ------------------------------------------------------------
	\begin{frame}{Little's Law Example}
		If:
		\begin{itemize}
			\item $\lambda = 50$ req/s
			\item $W = 0.2$ s (200ms)
		\end{itemize}
		

		\[
		L = 50 \times 0.2 = 10
		\]

		\begin{alertblock}{Interpretation}
			On average, 10 requests are inside the system at once.
		\end{alertblock}
	\end{frame}
	
	% ------------------------------------------------------------
	\begin{frame}{Student Exercise (1 minute)}
		Compute $L$ if:
		\[
		\lambda = 80~\text{req/s},\quad W = 0.25~\text{s}
		\]
		

		\[
		L = 80 \times 0.25 = 20
		\]
		

		\begin{block}{Meaning}
			Your system must handle about 20 concurrent requests on average.
		\end{block}
	\end{frame}
	
	% ============================================================
	\section{Context Switching Cost}
	% ============================================================
	
	% ------------------------------------------------------------
	\begin{frame}{Context Switching Cost --- Hidden Overhead}
		Threads are not free. Context switching can imply:
		\begin{itemize}
			\item saving/restoring registers
			\item pipeline disruption
			\item cache pollution
			\item scheduler overhead
		\end{itemize}
		

		\begin{alertblock}{When thread count is huge}
			The OS scheduler thrashes and p95 latency explodes.
		\end{alertblock}
	\end{frame}
	
	% ------------------------------------------------------------
	\begin{frame}{Context Switching Example}
		Assume:
		\begin{itemize}
			\item switch cost = $5~\mu s$
			\item $100{,}000$ switches/sec
		\end{itemize}
		

		\[
		100{,}000 \times 5\mu s = 0.5s
		\]
		

		\begin{alertblock}{Interpretation}
			Half the CPU time is wasted switching.
		\end{alertblock}
	\end{frame}
	
	% ============================================================
	\section{Synthesis and Activities}
	% ============================================================
	
	% ------------------------------------------------------------
	\begin{frame}{Synthesis}
		\begin{itemize}
			\item \textbf{Amdahl:} parallelism has a ceiling.
			\item \textbf{Little:} latency creates queues.
			\item \textbf{Context switching:} too many threads hurt.
		\end{itemize}
		

		\begin{block}{Balance}
			Idle time exploitation \quad vs \quad scheduling overhead.
		\end{block}
	\end{frame}
	
	% ------------------------------------------------------------
	\begin{frame}{Transition Activity: Predict -- Then -- Test}
		\textbf{Scenario:} ``We will add threads next.''
		
		Predict what happens to:
		\begin{enumerate}
			\item Latency (p50)
			\item Throughput
			\item CPU usage
		\end{enumerate}
		

		Write predictions now. We test in Part 2 experiments.
	\end{frame}
	
	% ------------------------------------------------------------
	\begin{frame}{Common Misconceptions Quiz (Pair Defense)}
		True/False (defend your answer in pairs):
		\begin{enumerate}
			\item High throughput always means low latency. (False)
			\item Adding more threads always improves performance. (False)
			\item Parallelism requires multiple cores; concurrency does not. (True)
			\item If CPU is 10\%, the system is underutilized. (Depends)
			\item Concurrency is only useful for I/O-bound tasks. (Mostly true; nuanced)
		\end{enumerate}
	\end{frame}
	
	% ------------------------------------------------------------
	\begin{frame}{Wrap + Homework Bridge}
		\textbf{Students repeat back:}
		\begin{itemize}
			\item Concurrency $\neq$ speed; concurrency hides waiting
			\item Latency $\neq$ throughput (both matter)
			\item Most servers are I/O-bound
		\end{itemize}
		

		\textbf{Homework (mandatory):}
		\begin{itemize}
			\item Run \texttt{AmdahlCalculator}
			\item Paste output table (or screenshot) into your report
			\item Answer: ``Why does speedup plateau?''
		\end{itemize}
	\end{frame}
	
	% ------------------------------------------------------------
	\begin{frame}{Exit Ticket (2 minutes)}
		Answer quickly:
		\begin{enumerate}
			\item Define latency in one sentence.
			\item Why can CPU be low while latency is high?
			\item True/False: more threads always help. (one keyword)
		\end{enumerate}
		

		\begin{alertblock}{Key takeaway}
			You cannot parallelize waiting away. You must hide it.
		\end{alertblock}
	\end{frame}
	
\end{document}
